\section{Свойства непрерывных на отрезке функций:
достижение точной верхней и нижней грани.
Теорема о промежуточных значениях. (Beta)}

\subsection{Ограниченность непрерывной функции на отрезке}

\begin{theorem}[Теорема Вейерштрасса]
Если функция $f$ непрерывна на отрезке $[a,b]$, то она ограничена на этом отрезке.
\end{theorem}

\begin{Proof}
Предположим противное: пусть функция не ограничена сверху.
Тогда существует последовательность $x_n \in [a,b]$, такая что
$f(x_n)\to +\infty$.

Так как отрезок $[a,b]$ компактен, из последовательности $x_n$
можно выделить сходящуюся подпоследовательность
$x_{n_k}\to x_0\in[a,b]$.

По непрерывности:
\[
f(x_{n_k}) \to f(x_0),
\]
что противоречит $f(x_{n_k})\to +\infty$.
Следовательно, функция ограничена.
\end{Proof}

\subsection{Достижение точной верхней и нижней грани}

\begin{theorem}[О достижении экстремумов]
Если функция $f$ непрерывна на отрезке $[a,b]$, то существуют точки
$x_{\min}, x_{\max} \in [a,b]$, такие что
\[
f(x_{\min})=\min_{[a,b]} f(x),
\qquad
f(x_{\max})=\max_{[a,b]} f(x).
\]
\end{theorem}

\begin{remark}
То есть точные верхняя и нижняя грани совпадают соответственно
с максимумом и минимумом функции и достигаются в некоторых точках.
\end{remark}

\subsection{Теорема о промежуточных значениях}

\begin{theorem}
Если функция $f$ непрерывна на $[a,b]$, то она принимает все значения
между $f(a)$ и $f(b)$.
\end{theorem}

То есть если число $C$ лежит между $f(a)$ и $f(b)$, то существует
точка $c\in[a,b]$, такая что
\[
f(c)=C.
\]

\begin{Proof}[Идея доказательства]
Рассматривается множество
\[
E=\{x\in[a,b]: f(x)\le C\}.
\]
Берётся его точная верхняя грань $c=\sup E$.
Используя непрерывность функции, доказывается, что $f(c)=C$.
\end{Proof}

\subsection{Следствие (Теорема Больцано)}

Если $f$ непрерывна на $[a,b]$ и
\[
f(a)\cdot f(b)<0,
\]
то существует $c\in(a,b)$ такая, что
\[
f(c)=0.
\]

\subsection{Геометрический смысл}

График функции, непрерывной на отрезке,
можно провести, не отрывая карандаша от бумаги.
Он не имеет разрывов и «скачков».

\subsection{Примеры}

\begin{example}
Рассмотрим функцию
\[
f(x)=x^3-x
\]
на отрезке $[-2,2]$.

Она непрерывна (многочлен), следовательно:
\begin{itemize}
\item ограничена;
\item достигает максимума и минимума;
\item принимает все значения между $f(-2)$ и $f(2)$.
\end{itemize}
\end{example}

\begin{example}
Докажем, что уравнение
\[
x^3+x-1=0
\]
имеет корень на $(0,1)$.

Пусть
\[
f(x)=x^3+x-1.
\]
Тогда
\[
f(0)=-1<0,
\qquad
f(1)=1>0.
\]
Функция непрерывна, значит существует
$c\in(0,1)$ такое, что $f(c)=0$.
\end{example}
